% Chapter 6 -- Configuring Your System: CONFIG.SYS and the Environment.
% Sources (Law 1): os/milton/config_sys.c, os/milton/sysinit.c (FILES cap
% 20; BREAK; DEVICE=ANSI.SYS; others recorded); spec/dos_config_sys_
% baseline.txt (standard settings); os/milton/env.c (512 bytes);
% os/milton/command.c (builtin_set, builtin_prompt, command_repl COMSPEC /
% PROMPT / PATH, run_external PATH search).  PATH as a command:
% initech-sy5f.
% GATES: test-config-sys test-sysinit test-env test-command test-pubs-screens
\Gates{test-config-sys, test-sysinit, test-env, test-command, test-pubs-screens}
\chapter{Configuring Your System}

\section{The CONFIG.SYS File}
\index{CONFIG.SYS@\texttt{CONFIG.SYS}}
When InitechDOS starts, it reads the file \file{CONFIG.SYS} from the root
directory of drive~A.  This file holds \term{configuration commands}, one to
a line, in the form \ph{command}\kw{=}\ph{value}.  If there is no
\file{CONFIG.SYS}, InitechDOS uses its standard settings.  Most users never
need to change them.

These are the configuration commands; the \doctitle{InitechDOS Reference},
Chapter~3, describes each one:
\par\medskip
\begin{tabular}{@{}l@{\quad}l@{\quad}p{44mm}@{}}
{\imSideHeadFont Command} & {\imSideHeadFont Standard} & {\imSideHeadFont What it sets}\\[2pt]
\cmd{BREAK}     & \texttt{ON}      & Whether \key{Ctrl-C} is checked often\\
\cmd{BUFFERS}   & \texttt{20}      & The number of disk buffers\\
\cmd{DEVICE}    & \texttt{ANSI.SYS} & A device driver to install\\
\cmd{FILES}     & \texttt{20}      & How many files can be open at once\\
\cmd{INSTALL}   & \texttt{SHARE.EXE} & A program to load at start-up\\
\cmd{LASTDRIVE} & \texttt{Z}       & The last drive letter\\
\cmd{SHELL}     & \texttt{COMMAND.COM} & The command processor\\
\end{tabular}
\par\medskip
In this version of InitechDOS, \cmd{FILES}, \cmd{BREAK} and
\texttt{DEVICE=ANSI.SYS} take effect; the other commands are checked and
recorded, and do not change the way the system runs.  To create a
\file{CONFIG.SYS} file, use \cmd{COPY CON} as described in Chapter~5, for
example with the single line \texttt{BREAK=OFF}, and then restart the
system.  Lines InitechDOS does not recognize are ignored.

\section{The Environment}
\index{environment}\index{SET@\texttt{SET}}
The \term{environment} is a small area of memory in which InitechDOS keeps
named settings, called \term{environment variables}, for itself and for the
programs you run.  The \cmd{SET} command lists them and changes them.  In
this exercise you add a variable of your own:
\begin{Steps}
\item Type \type{CLS} and press \key{Enter}.
\item Type \type{SET DEPT=PUBS} and press \key{Enter}.
\item Type \type{SET} and press \key{Enter}.
\end{Steps}

\Screen[rows=7, label=fig:set]{screens/set.txt}{The environment}

InitechDOS sets three variables itself:

\SideHead{COMSPEC}Where the command processor is.

\SideHead{PROMPT}The form of the command prompt (see below).

\SideHead{PATH}The directories InitechDOS searches for a program or batch
file that is not in the current directory, separated by semicolons.

To remove a variable, type its name and an equal sign with nothing after
it: \type{SET DEPT=}.  The environment holds 512 characters; when it is
full, \cmd{SET} displays \texttt{Out of environment space}.  Batch files can
use variables by name: \texttt{\%DEPT\%} in a batch file stands for
\texttt{PUBS}.

\section{The Search Path}
\index{PATH@\texttt{PATH}}\index{search path}
When you type the name of a program that is not in the current directory,
InitechDOS looks in each directory named in \texttt{PATH}, in order.  The
practice diskette's \file{AUTOEXEC.BAT} sets it with
\type{SET PATH=A:\bs APPS}, so that the utility programs in \file{\bs APPS}
can be run from any directory.  To search more than one directory, separate
them with semicolons: \type{SET PATH=A:\bs APPS;A:\bs SAMPLES}.

The \cmd{PATH} command does the same with less typing:
\type{PATH A:\bs APPS}.  Typed alone, \cmd{PATH} displays the search path;
\type{PATH ;} clears it, so that InitechDOS looks only in the current
directory, and \cmd{PATH} then displays \texttt{No Path}.
\begin{Steps}
\item Type \type{CLS} and press \key{Enter}.
\item Type \type{PATH} and press \key{Enter} to see the search path that
  \file{AUTOEXEC.BAT} set.
\item Type \type{PATH ;} and press \key{Enter}, then \type{PATH} again.
\item Type \type{PATH A:\bs APPS} and press \key{Enter}, then \type{PATH}
  again to see the search path restored.
\end{Steps}

\Screen[rows=9, label=fig:path]{screens/path.txt}{Displaying, clearing and
  setting the search path}

\section{Changing the Command Prompt}
\index{PROMPT@\texttt{PROMPT}}
The \cmd{PROMPT} command changes the command prompt.  A dollar sign and a
letter stand for information: \kw{\$P} is the current drive and directory,
\kw{\$G} the \texttt{>} character, \kw{\$N} the drive letter, and \kw{\$T}
the time.  Any other text is displayed as you type it.  The standard prompt
is \kw{\$P\$G}.
\begin{Steps}
\item Type \type{CLS} and press \key{Enter}.
\item Type \type{PROMPT \$N\$G} and press \key{Enter}.  The prompt becomes
  \texttt{A>}.
\item Type \type{PROMPT} and press \key{Enter}.  The standard prompt comes
  back.
\end{Steps}

\Screen[rows=3, label=fig:prompt]{screens/prompt.txt}{Changing the prompt and back}

The \doctitle{InitechDOS Reference} lists all the \kw{\$} codes.  We
recommend a prompt that includes \kw{\$P}, so that you always know which
directory you are in.
